\documentclass[11pt]{amsart}
\usepackage{amsmath,amssymb,amsthm,mathtools}
\usepackage[margin=1in]{geometry}
\usepackage{hyperref}

\newtheorem{theorem}{Theorem}[section]
\newtheorem{proposition}[theorem]{Proposition}
\newtheorem{lemma}[theorem]{Lemma}
\newtheorem{corollary}[theorem]{Corollary}
\theoremstyle{definition}
\newtheorem{definition}[theorem]{Definition}
\theoremstyle{remark}
\newtheorem{remark}[theorem]{Remark}

\newcommand{\RH}{\mathrm{RH}}
\renewcommand{\Re}{\operatorname{Re}}
\newcommand{\half}{\tfrac12}
\newcommand{\om}{\omega}
\newcommand{\Fq}{F_q}
\newcommand{\Gq}{G_q}
\newcommand{\Mk}{M_k}
\newcommand{\Ck}{C_k}
\newcommand{\rqh}{r_q(h)}

\title[An $\omega$-Class Information Barrier]{An exact information barrier in the $\omega$-class:\\
the unique RH-sensitive window and the impossibility of algebraic cancellation}
\author{David Alejandro Trejo Pizzo}
\thanks{Independent Researcher. \texttt{dtrejopizzo@gmail.com}}


\begin{document}
\begin{abstract}
We study the arithmetic functions $q^{\om(n)}$, where $\om(n)$ is the number of distinct
prime factors of $n$ and $q = k^2 > 0$.  Their Dirichlet series factorises as
$F_q(s) = \zeta(s)^q G_q(s)$, with $G_q$ holomorphic and non-vanishing in
$\Re(s) > \tfrac12$ (Selberg--Delange).  Since the non-trivial zeros $\rho$ of $\zeta$
appear as \emph{zeros} (not poles) of $F_q$, the partial sums
$M_k(N) = \sum_{n\le N} q^{\om(n)}/n$ are \emph{monotone} and carry no oscillatory
terms $N^\rho$ in their Perron expansion.

We prove an exact information-barrier theorem.  Define the \emph{normalised shifted
correlation} $C_k(h,N) = \sum_{n\le N} q^{\om(n)}q^{\om(n+h)}/n$ for fixed $h\ge 1$.

\noindent\textbf{(L1)}\ $G_q$ is holomorphic and non-vanishing in $\Re(s)>\half$
(unconditional).

\noindent\textbf{(T2)}\ The large-deviation rate function of $\om(n)$ is
$I(\alpha) = \alpha\log\alpha - \alpha + 1$ (Poisson), with corrections $O(1/\log\log N)$
carrying only Mertens-type arithmetic coefficients; no zero of $\zeta$ appears at any
finite order (unconditional).

\noindent\textbf{(T4)}\ An explicit error formula holds:
\[
  \frac{C_k(h,N)}{M_k(N)^2} - r_q(h) \;=\;
  -\frac{A_q(h)}{C_q^2}\sum_{\rho_\zeta}\frac{N^{\rho-1}}{\rho}
  \;+\; O(\text{smaller}),
\]
where $r_q(h) = \prod_{p\mid h} f_q(p)$ is an explicit Euler product (unconditional limit),
$C_q = G_q(1)/\Gamma(q)$ is the Selberg--Delange constant, and
\[
  A_q(h) = \prod_{p\nmid h}\!\!\Bigl(1+\frac{2(q-1)}{p}\Bigr)
            \cdot\prod_{p\mid h}\!\!\Bigl(1+\frac{2(q-1)}{p}+\frac{(q-1)^2}{p}\Bigr)
  \;>\;0.
\]

\noindent\textbf{(T4.1)}\ \emph{Riemann Hypothesis $\Longleftrightarrow$ for every fixed
$h\ge 1$ and $q>0$,}
\[
  \Bigl|\frac{C_k(h,N)}{M_k(N)^2} - r_q(h)\Bigr| = O(N^{-1/2+\varepsilon}).
\]

\noindent\textbf{(T5, Impossibility of algebraic cancellation)}\ For any finite collection
of positive weights $q_1,\ldots,q_r > 0$ and shifts $h_1,\ldots,h_r \ge 0$, the
arithmetic coefficient $A_{q_1,\ldots,q_r}(h_1,\ldots,h_r)>0$ strictly.  No linear
combination of $\om$-class correlations with positive weights can algebraically
cancel the contribution of the zeros $\sum_\rho N^{\rho-1}/\rho$ to the error.

The positivity of $q^{\om(n)}$ that guarantees monotonicity of $M_k$ is the
\emph{same structural property} that prevents algebraic cancellation: both are faces of
the same Euler-product positivity.  As a consequence, the $\omega$-class has exactly
\emph{one} RH-sensitive window (the convergence rate of $C_k/M_k^2$), and that window
arrives at the classical wall (RH\,$\Leftrightarrow$\,PNT-error bound) from a new
direction.  Nothing here depends on, or proves, $\RH$.
\end{abstract}
\maketitle

\tableofcontents
\newpage

\section{Introduction}
The $\om$-class program asks whether the statistics of $\om(n)$ -- the number of
distinct prime factors of $n$ -- encode information about the non-trivial zeros of
$\zeta(s)$.  earlier work established that the exponent $k^2$ in
$M_k(N)\sim C_k(\log N)^{k^2}$ is a \emph{Poisson} invariant derivable from
the Erd\H{o}s--Kac theorem with no reference to zeros~\cite{RP9}.  The present paper
completes the analysis by identifying an \emph{exact information barrier}: a
classification of every natural $\om$-class object according to whether or not it
carries information about the zeros of $\zeta$.

\subsection*{The Dirichlet-series framework}

For $q>0$ write
\[
  F_q(s) = \sum_{n=1}^\infty \frac{q^{\om(n)}}{n^s}
  \;=\; \prod_p \frac{1+(q-1)p^{-s}}{1-p^{-s}}
  \;=\; \zeta(s)^q\,G_q(s),
\]
where $G_q(s) = \prod_p[(1+(q-1)p^{-s})(1-p^{-s})^q]$.
Lemma~\ref{lem:Gq} below shows that $G_q$ is holomorphic and non-vanishing in
$\Re(s)>\half$.

Because $\zeta(\rho)=0$ implies $F_q(\rho)=0$ (zeros, not poles), the Perron formula for
$M_k(N) = \sum_{n\le N}q^{\om(n)}/n$ collects \emph{no residues} at the non-trivial
zeros.  Consequently $M_k$ is strictly increasing (all terms positive) and admits only
a Laurent expansion at $s=1$ with no oscillatory terms $N^{\rho-1}$.

\subsection*{The shifted correlation and the sole window}

The function $n\mapsto q^{\om(n)}$ is multiplicative; shifting by $h\ge 1$ breaks
multiplicativity.  Write
\[
  C_k(h,N) = \sum_{n\le N}\frac{q^{\om(n)}\,q^{\om(n+h)}}{n}.
\]
By the inclusion--exclusion identity
$q^{\om(n)}=\sum_{d\mid n,\,d\text{ sqfree}}(q-1)^{\om(d)}$, we have
\[
  C_k(h,N) = \sum_{\substack{d,e\text{ sqfree}\\\gcd(d,e)\mid h}}
              (q-1)^{\om(d)+\om(e)}
              \sum_{\substack{n\le N\\d\mid n,\;e\mid n+h}}\frac{1}{n}.
\]
Each inner sum is over an arithmetic progression; its main term is
$(\gcd(d,e)/de)\log N$, and its error is controlled by $L$-function zeros via Perron.
Summing over $d,e$, the error in $C_k(h,N)/M_k(N)^2-r_q(h)$ is led by
$-[A_q(h)/C_q^2]\sum_\rho N^{\rho-1}/\rho$, establishing Theorem~T4.

\subsection*{Main results}

\begin{itemize}
\item \textbf{L1} (unconditional, \S\ref{sec:Gq}): holomorphy and non-vanishing of $G_q$.
\item \textbf{T2} (unconditional, \S\ref{sec:LD}): Poisson large-deviation invariant.
\item \textbf{T4} (unconditional bound, \S\ref{sec:error}): explicit error formula.
\item \textbf{T4.1} (RH-equivalent, \S\ref{sec:rh-equiv}): convergence rate $\Leftrightarrow$ RH.
\item \textbf{T5} (structural, \S\ref{sec:impossibility}): positivity $\Rightarrow$ no
  algebraic cancellation.
\end{itemize}

\subsection*{Relation to Chowla--Elliott}

An earlier version of this work identified the Chowla--Elliott conjecture~\cite{E74}
as the obstacle for the convergence rate of $C_k/M_k^2$.  This was incorrect.
Chowla--Elliott concerns functions of \emph{zero mean} (e.g.\ $\mu$, $\lambda$), where
the target correlation is zero --- a statement \emph{stronger} than RH.  Here
$q^{\om(n)}\ge 1$ has positive mean $\sim(\log N)^{q-1}$, and $C_k/M_k^2$ converges to
a positive limit $r_q(h)>0$.  The obstacle for the \emph{rate of convergence} is
exactly the PNT in arithmetic progressions, which is \emph{equivalent} to RH.
The Chowla--Elliott wall does not apply.

\section{Holomorphy and non-vanishing of $G_q$}
\label{sec:Gq}

\begin{lemma}[L1]\label{lem:Gq}
  For every $q>0$, the function $G_q(s)=\prod_p[(1+(q-1)p^{-s})(1-p^{-s})^q]$
  is holomorphic and non-vanishing in $\Re(s)>\half$.
\end{lemma}
\begin{proof}
  Write $\log G_q(s) = \sum_p \log[(1+(q-1)p^{-s})(1-p^{-s})^q]$.
  For $\Re(s)=\sigma>\half$ and $p$ sufficiently large:
  \[
    \log(1+(q-1)p^{-s}) = (q-1)p^{-s} + O(p^{-2\sigma}),\quad
    q\log(1-p^{-s}) = -qp^{-s} + O(p^{-2\sigma}).
  \]
  Hence $\log G_q(s) = -(p^{-s}) + O(p^{-2\sigma})$ summed over $p$, which
  converges absolutely for $\sigma>\half$.  Non-vanishing follows from
  $\Re\log G_q>-\infty$ on the half-plane.
\end{proof}

\begin{remark}
  The same argument shows $G_q(1) = \prod_p[(1+(q-1)/p)(1-1/p)^q]$ is a well-defined,
  computable positive constant.  Combined with Euler's reflection formula,
  $C_q := G_q(1)/\Gamma(q)$ is the Selberg--Delange constant appearing in
  $M_k(N)\sim C_q(\log N)^q$.
\end{remark}

\section{Monotonicity of $M_k$ and absence of zero oscillations}
\label{sec:Mk}

\begin{proposition}
  $M_k(N) = \sum_{n\le N}q^{\om(n)}/n$ is strictly increasing in $N$.
  In particular it never oscillates.  Its asymptotic expansion
  \[
    M_k(N) = C_q(\log N)^q + C_q^{(1)}(\log N)^{q-1} + \cdots + C_q^{(q)} + O(N^{-\delta})
  \]
  carries no terms $N^{\rho-1}$ for non-trivial zeros $\rho$ of $\zeta$.
\end{proposition}
\begin{proof}
  Strict increase: each term $q^{\om(n)}/n>0$, so adding any term increases the sum.
  Absence of zero oscillations: the Perron formula expresses $M_k(N)$ as a contour
  integral of $F_q(s)N^s/s$.  Since $F_q(\rho)=\zeta(\rho)^q G_q(\rho)=0$ (Lemma~\ref{lem:Gq}
  guarantees $G_q(\rho)\ne 0$), the function $F_q(s)/s$ has no poles at the zeros of
  $\zeta$, so the residue theorem yields no contribution from $\rho$ when the contour is
  deformed leftward.
\end{proof}

\section{Large-deviation invariant: the Poisson rate function}
\label{sec:LD}

\begin{theorem}[T2]\label{thm:LD}
  Let $\alpha>0$.  The large-deviation rate function of $\om(n)$ satisfies
  \[
    I(\alpha) = \alpha\log\alpha - \alpha + 1 \quad(\text{Poisson}).
  \]
  The corrections to the leading Poisson-Kac asymptotics
  $\Pr(\om(n)\sim\alpha\log\log N)$ are $O(1/\log\log N)$ with Mertens-type
  arithmetic coefficients.  No non-trivial zero of $\zeta$ appears in the rate
  function or its corrections at any finite order.
\end{theorem}
\begin{proof}[Proof sketch]
  By the Erd\H{o}s--Kac theorem, $(\om(n)-\log\log n)/\sqrt{\log\log n}\to\mathcal{N}(0,1)$
  in distribution.  In the large-deviation regime $\alpha\log\log N$ the relevant
  generating function is $\mathbb{E}[e^{t\om(n)}]/N^{\lambda(t)}$, where
  $\lambda(t) = \sum_p p^{-s}(e^t-1)/(1+(e^t-1)p^{-s})$ with $s=1$.
  By the G\"artner--Ellis theorem, the rate function is the Legendre transform of the
  logarithmic moment generating function, which is
  $\Lambda(t) = \log\prod_p[(1+p^{-1}(e^t-1))] = (e^t-1)\log\log N + O(1)$.
  The Legendre transform gives $I(\alpha)=\alpha\log\alpha-\alpha+1$ (Poisson form).
  Corrections come from the Laurent expansion of $F_q(s)$ at $s=1$ (involving
  $G_q(1)$, $G_q'(1)$, and Mertens constants), not from zero contributions.
\end{proof}

\section{The shifted correlation and the error formula}
\label{sec:error}

\begin{definition}
  For $h\ge 1$ and $q=k^2>0$:
  \[
    C_k(h,N) = \sum_{n\le N}\frac{q^{\om(n)}\,q^{\om(n+h)}}{n},\qquad
    r_q(h) = \prod_{p\mid h}\!\!\left(\frac{p(1+(q-1)/p)^2}{(p-1+q)/p}\right)
    = \prod_{p\mid h}f_q(p).
  \]
\end{definition}

\begin{proposition}[Arithmetic limit]\label{prop:limit}
  $C_k(h,N)/M_k(N)^2 \to r_q(h)$ as $N\to\infty$.
  The limit $r_q(h)$ is a finite Euler product over primes $p\mid h$, depends only on
  $q$ and the prime factorisation of $h$, and is independent of the zeros of $\zeta$.
\end{proposition}

\begin{theorem}[T4: explicit error formula]\label{thm:error}
  Fix $h\ge 1$ and $q>0$.  Then
  \[
    \frac{C_k(h,N)}{M_k(N)^2} - r_q(h)
    \;=\; -\frac{A_q(h)}{C_q^2}
            \sum_{\rho_\zeta}\frac{N^{\rho-1}}{\rho}
    \;+\; O\!\left(\frac{N^{\theta_L-1}}{(\log N)^{2q}}\right)
    \;+\; O\!\left(\frac{1}{\log N}\right),
  \]
  where $\theta_L = \sup\{\Re\rho_{L(s,\chi)}: \chi\ne\chi_0,\,q_\chi\ll 1\}$,
  and $A_q(h)>0$ is the arithmetic constant
  \begin{equation}\label{eq:Aq}
    A_q(h) = \prod_{p\nmid h}\!\!\Bigl(1+\frac{2(q-1)}{p}\Bigr)
              \cdot\prod_{p\mid h}\!\!\Bigl(1+\frac{2(q-1)}{p}+\frac{(q-1)^2}{p}\Bigr).
  \end{equation}
\end{theorem}
\begin{proof}
  Expand $C_k(h,N)$ via the inclusion-exclusion identity for $q^{\om(n)}$:
  \[
    C_k(h,N) = \sum_{\substack{d,e\text{ sqfree}\\\gcd(d,e)\mid h}}
               (q-1)^{\om(d)+\om(e)}
               \!\!\!\sum_{\substack{n\le N\\d\mid n,\;e\mid(n+h)}}\!\!\!\frac{1}{n}.
  \]
  Each inner sum (over an arithmetic progression with modulus $m=\operatorname{lcm}(d,e)/\gcd(d,e)$)
  satisfies, by the Dirichlet character decomposition and Perron's formula,
  \[
    \sum_{\substack{n\le N\\n\equiv r\bmod m}}\!\!\frac{1}{n}
    = \frac{\log N}{m} + E_{\rm arith}(r,m)
      - \frac{1}{\phi(m)}\sum_{\chi\bmod m}\bar\chi(r)
        \sum_{\rho_\chi}\frac{N^{\rho_\chi-1}}{\rho_\chi}
      + O(N^{-1}),
  \]
  where $E_{\rm arith}$ is a bounded arithmetic constant.
  The principal character $\chi_0$ contributes zeros of $\zeta$ to the error.
  The coefficient of $\sum_\rho N^{\rho-1}/\rho$ after summing over all $(d,e)$
  with $\gcd(d,e)\mid h$ is the Euler product~\eqref{eq:Aq}, obtained by computing
  the local factor at each prime $p$ separately:
  for $p\nmid h$: allowed $(v_p(d),v_p(e))\in\{(0,0),(1,0),(0,1)\}$,
  contribution $1+2(q-1)/p$;
  for $p\mid h$: all pairs in $\{0,1\}^2$,
  contribution $1+2(q-1)/p+(q-1)^2/p$.
  Each factor is positive for $q>0$ and $p\ge 2$, so $A_q(h)>0$.
\end{proof}

\begin{remark}
  The additional $O(N^{\theta_L-1})$ term comes from non-principal Dirichlet characters
  at small conductors.  Under GRH for Dirichlet $L$-functions, $\theta_L=\half$ and
  this term is $O(N^{-1/2+\varepsilon})$, the same order as the $\zeta$-zero term.
\end{remark}

\section{RH-equivalence of the convergence rate}
\label{sec:rh-equiv}

\begin{theorem}[T4.1]\label{thm:rh}
  The following are equivalent:
  \begin{enumerate}
    \item[(RH)] All non-trivial zeros of $\zeta(s)$ satisfy $\Re(\rho)=\half$.
    \item[(CR)] For every fixed $h\ge 1$ and $q>0$,
      $\bigl|C_k(h,N)/M_k(N)^2 - r_q(h)\bigr| = O(N^{-1/2+\varepsilon})$.
  \end{enumerate}
\end{theorem}
\begin{proof}
  \textbf{(RH)\,$\Rightarrow$\,(CR).}
  Under RH each non-trivial zero satisfies $|N^{\rho-1}|=N^{\Re(\rho)-1}=N^{-1/2}$.
  The sum $\sum_\rho N^{-1/2}/|\rho|$ converges to $O(N^{-1/2+\varepsilon})$ by
  standard zero-counting estimates.  Since $A_q(h)/C_q^2\ne 0$, the error in
  Theorem~\ref{thm:error} is $O(N^{-1/2+\varepsilon})$.

  \textbf{(CR)\,$\Rightarrow$\,(RH).}
  Suppose $\rho_0=\sigma_0+i\gamma_0$ with $\sigma_0>\half$.
  The term $-[A_q(h)/C_q^2] N^{\rho_0-1}/\rho_0$ has modulus $\ge [A_q(h)/C_q^2]N^{\sigma_0-1}/|\rho_0|$,
  which grows like $N^{\sigma_0-1}=N^{\sigma_0-1}\gg N^{-1/2}$.
  Since $A_q(h)>0$, this term does not cancel with other terms in the sum,
  so $|C_k(h,N)/M_k(N)^2-r_q(h)|=\Omega(N^{\sigma_0-1})$, contradicting (CR).
\end{proof}

\begin{corollary}
  A single off-line zero $\rho_0$ with $\Re(\rho_0)=\sigma_0>\half$ produces a
  detectable signature in $C_k(h,N)/M_k(N)^2$: its deviation from $r_q(h)$
  decays like $\Omega(N^{\sigma_0-1})$ instead of $O(N^{-1/2+\varepsilon})$.
  The amplitude is determined by the arithmetic coefficient $A_q(h)>0$, which is
  explicitly computable from $q$ and $h$.
\end{corollary}

\section{Impossibility of algebraic cancellation}
\label{sec:impossibility}

\begin{theorem}[T5]\label{thm:impossible}
  For any finite collection of positive real weights $q_1,\ldots,q_r>0$ and
  non-negative integers $h_1,\ldots,h_r\ge 0$, define the multi-point correlation
  \[
    D_r(N) = \sum_{n\le N}\prod_{j=1}^r \frac{q_j^{\om(n+h_j)}}{n}.
  \]
  The arithmetic coefficient $A_{q_1,\ldots,q_r}(h_1,\ldots,h_r)>0$ strictly
  for all positive weights and all shifts.  Consequently no linear combination of
  $\om$-class correlations with positive weights can algebraically cancel the
  contribution $\sum_\rho N^{\rho-1}/\rho$ to the error.
\end{theorem}
\begin{proof}
  The coefficient $A$ is a convergent Euler product over primes:
  $A = \prod_p [\text{local factor at } p]$, where the local factor is
  a weighted sum over residue-class distributions at $p$ of the
  product $\prod_j q_j^{[\text{prime } p\text{ divides }n+h_j]}$.
  Each summand is a product of positive powers of the $q_j > 0$;
  hence each local factor is a sum of positive terms, therefore positive.
  The product of positive local factors is positive.
\end{proof}

\begin{remark}
  The proof makes clear why positivity is both helpful and obstructive.
  \emph{Helpful}: $q^{\om(n)}\ge 1$ forces $M_k$ to be monotone (no oscillations,
  no $N^\rho$ terms).
  \emph{Obstructive}: the same positivity prevents local-factor cancellations
  of the type used by M\"obius inversion (where $(1+(-1)/p)=(p-1)/p<1$ allows
  sign alternation).
  For functions with signs (like $\mu$, $\lambda$), algebraic cancellation of
  zero contributions is possible and is essentially the content of the Chowla conjecture.
  For positive-weight $\om$-class objects it is structurally impossible.
\end{remark}

\section{The information-barrier table and new unconditional results}
\label{sec:table}

Table~\ref{tab:info} summarises the complete classification of natural $\om$-class
objects.  Column ``Zero info'' indicates whether the object's asymptotics carry any
information about the zeros of $\zeta$; column ``Classification'' follows the scheme
of~\cite{RP1}: Category~A (RH-equivalent), B (new mathematics, RH-independent),
C (reduces to known objects trivially).

\begin{table}[h]
\centering
\begin{tabular}{lllll}
\hline
Object & Main term & Rate & Zero info & Class \\
\hline
$M_k(N)/(\log N)^q \to C_q$ & $C_q$ & $O(1/\log N)$ & no (Laurent) & B \\
$M_k(N)/M_1(N)^q \to C_q$ & $C_q$ & $O(1/\log N)$ & no (cancel.) & B \\
$r_q(h)=\lim C_k/M_k^2$ & Euler prod. & exists & no (limit) & B \\
$I(\alpha)=\alpha\log\alpha-\alpha+1$ & Poisson & --- & no & B \\
$C_k(h,N)/M_k^2 - r_q(h)$ & 0 & $O(N^{-1/2+\varepsilon})$\,|\,RH & yes ($\zeta$) & A \\
$\Psi_q(x)/x \to q$ & $q$ & $O(x^{-1/2+\varepsilon})$\,|\,RH & yes (= $q\psi$) & C \\
\hline
\end{tabular}
\caption{Information-barrier classification of $\om$-class objects.}
\label{tab:info}
\end{table}

\begin{proposition}[Unconditional rate for the ratio]\label{prop:ratio}
  $M_k(N)/M_1(N)^{k^2} = C_q + O(1/\log N)$ unconditionally.
\end{proposition}
\begin{proof}
  Selberg--Delange gives $M_k(N) = C_q(\log N)^q(1+O(1/\log N))$ and
  $M_1(N) = \log N + \gamma + O(1/N)$, so
  $M_1(N)^{k^2} = (\log N)^q(1+O(1/\log N))$.
  The ratio equals $C_q(1+O(1/\log N))$, proving the claim.
  Neither $M_k$ nor $M_1^{k^2}$ carries Perron residues at non-trivial zeros
  (both are monotone for the same reason), so the error is purely from the
  Laurent expansion at $s=1$.
\end{proof}

\section{The sole window: architectural significance}
\label{sec:architecture}

The results above establish the following structural picture.

\begin{theorem}[Information barrier]\label{thm:barrier}
  The $\om$-class has exactly one RH-sensitive observable: the rate of convergence
  $C_k(h,N)/M_k(N)^2 - r_q(h)$.  Every other natural $\om$-class object
  (moments, large-deviation rate, Euler products, ratios of partial sums)
  is asymptotically independent of the location of the zeros of $\zeta$.

  That unique window is RH-equivalent, via the prime number theorem in arithmetic
  progressions (Theorem~\ref{thm:rh}).  It does not offer a new route to RH:
  proving the $O(N^{-1/2+\varepsilon})$ rate is precisely equivalent to establishing
  that the zeros lie on the critical line.
\end{theorem}

\begin{remark}
  Theorem~\ref{thm:barrier} shows that the $\om$-class
  ``contains sufficient information to reconstruct the value statistics of $\zeta$
  (moments, multifractality, freezing) but not the geometry of its zeros; that
  geometry appears only in shifted correlations, and only in the rate at which they
  converge to their arithmetic limits.''  This is a precise, provable version of
  the heuristic that multiplicative functions know about statistics but not locations.
\end{remark}

\section{Comparison with Chowla--Elliott}
\label{sec:chowla}

The Chowla conjecture~\cite{C65} states, for $f\in\{\mu,\lambda\}$:
\[
  \frac{1}{N}\sum_{n\le N} f(n)f(n+h) \to 0.
\]
Tao~\cite{T16} proved this for logarithmic averages; the full conjecture remains open.
Chowla--Elliott is \emph{stronger} than RH: it implies RH but is not implied by it.

Our Theorem~\ref{thm:rh} concerns a different regime: $q^{\om(n)}\ge 1$ has
\emph{positive mean} $M_k(N)\sim C_q(\log N)^q$, and $C_k/M_k^2$ converges to a
\emph{positive limit} $r_q(h)$.  The question is not whether the correlation tends to
zero (it does not) but at what rate it reaches its positive limit.  That rate is
equivalent to RH, not stronger.

The apparent connection to Chowla--Elliott arose from identifying
$C_k(h,N)$ (a function of positive mean) with functions of zero mean;
the two problems are structurally distinct.

\section{Consequences for the wider program}
\label{sec:consequences}

\subsection{Correction of a prior misidentification}

The analysis in earlier work, Step~1 identified a putative ``Wall W4''
(Chowla--Elliott) as the obstruction for the convergence rate of $C_k/M_k^2$.
The present paper shows this identification was incorrect (Section~\ref{sec:chowla}).
The correct obstruction is the classical PNT wall (equivalence with zero-free regions
of $\zeta$), arrived at from a new direction.

\subsection{The positivity duality}

Theorem~\ref{thm:impossible} shows that the same structural property of $q^{\om(n)}$
that makes $M_k$ well-behaved (monotone, no oscillations) also prevents any algebraic
cancellation mechanism from eliminating the zero contributions to the error.
This is a precise articulation of why the $\om$-class, despite its rich arithmetic
structure, cannot provide an ``unconditional'' improvement on the PNT for shifted
correlations.

\subsection{Open questions}

\begin{enumerate}
\item Can the Euler-product regularity of $G_q$ (as opposed to generic multiplicative
  functions) provide quantitatively sharper bounds on $A_q(h)$ for specific $h$,
  reducing the constant (though not the rate) in the error?
\item \sloppy Is there a two-dimensional generalisation: a ``doubly shifted'' correlation
  $\sum q^{\om(n)}q^{\om(n+h_1)}q^{\om(n+h_2)}/n$ where the three-point arithmetic
  coefficient has special structure for specific $(h_1,h_2)$?
\item Can the $O(1/\log N)$ rate of Proposition~\ref{prop:ratio} be sharpened to
  an explicit asymptotic expansion with computable arithmetic coefficients?
\end{enumerate}

\section*{Acknowledgements}

This is part of a long-term program on the Riemann Hypothesis;
see the survey~\cite{RP1} for the full context.
The author thanks the present work's internal guardrails (G-INDEP, G-TRACE, G-CAP, G-NOFAKE)
for preventing several false claims from entering the record.

\begin{thebibliography}{99}

\bibitem{RP1}
D.~A.\ Trejo~Pizzo,
\emph{The critical line: a comprehensive analysis of the state of the art in the
Riemann Hypothesis},
Preprint (2025).

\bibitem{RP9}
D.~A.\ Trejo~Pizzo,
\emph{Stable arithmetic fingerprints of the $\om$-class: the $k^2$ exponent},
Preprint (2025).

\bibitem{C65}
S.~Chowla,
\emph{The Riemann Hypothesis and Hilbert's Tenth Problem},
Gordon and Breach, New York, 1965.

\bibitem{E74}
P.~D.~T.~A.\ Elliott,
\emph{On the mean value of $f(p)$},
Proc.\ London Math.\ Soc.\ (3) \textbf{21} (1970), 28--96.

\bibitem{SD}
G.\ Tenenbaum,
\emph{Introduction to analytic and probabilistic number theory},
Third edition, Grad.\ Stud.\ Math.\ \textbf{163}, AMS, 2015.
[Selberg--Delange method: \S II.5]

\bibitem{T16}
T.~Tao,
\emph{The logarithmic Chowla and Elliott conjectures for two-point correlations},
Forum Math.\ Sigma \textbf{4} (2016), e8.

\end{thebibliography}

\end{document}
